\documentclass[preview]{standalone}

\usepackage{amsmath}
\usepackage{amssymb}
\usepackage{stellar}
\usepackage{definitions}

\begin{document}

\id{psicologia-memoria-breve-lavoro}
\genpage

\section{Memoria sensoriale, breve termine e lavoro}

\begin{snippet}{memoria-temporanea-introduzione}
    Molte informazioni incontrate nel corso dell'esperienza non rimangono
    conservate stabilmente. Sono disponibili solo per brevi intervalli di tempo,
    attraverso sistemi come la memoria iconica, la memoria a breve termine e la
    memoria di lavoro.
\end{snippet}

\begin{snippetdefinition}{memoria-iconica-definition}{Memoria iconica}
    La \emph{memoria iconica}, o \emph{memoria sensoriale visiva}, è il sistema
    di memoria relativo al dominio visivo che permette di immagazzinare una
    grande quantità di informazioni per un intervallo di tempo molto breve.
\end{snippetdefinition}

\begin{snippetexample}{esperimento-sperling-example}{Esperimento di Sperling}
    George Sperling presentò ai partecipanti una matrice di lettere per un tempo
    brevissimo. Nella procedura di \emph{recupero totale}, i partecipanti
    dovevano elencare tutte le lettere ricordate. Nella procedura di
    \emph{recupero parziale}, invece, un segnale acustico alto, medio o basso
    indicava quale riga riportare. Poiché i partecipanti riuscivano a riportare
    accuratamente qualunque riga fosse indicata subito dopo la presentazione,
    Sperling concluse che l'intera matrice era stata acquisita dalla memoria
    iconica, ma svaniva molto rapidamente.
\end{snippetexample}

\includesnpt[width=45\%,src=/snippet/static-psychology/sperling-iconic-memory-matrix.jpg]{centered-img}

\begin{snippetnote}{memoria-iconica-immagine-eidetica-note}{Memoria iconica e immagine eidetica}
    La memoria iconica non coincide con la memoria fotografica. Il termine
    tecnico per la memoria fotografica è \emph{immagine eidetica}: chi possiede
    immagini eidetiche riesce a ricordare dettagli di un'immagine per un tempo
    molto più lungo rispetto alla memoria iconica.
\end{snippetnote}

\begin{snippetdefinition}{memoria-breve-termine-definition}{Memoria a breve termine}
    La \emph{memoria a breve termine}, o \emph{MBT}, consente di focalizzare le
    risorse cognitive su un numero limitato di rappresentazioni mentali per un
    periodo temporaneo.
\end{snippetdefinition}

\begin{snippetdefinition}{span-memoria-definition}{Span di memoria}
    Lo \emph{span di memoria} è la quantità di elementi che una persona riesce a
    mantenere temporaneamente in memoria e a riprodurre nell'ordine corretto.
\end{snippetdefinition}

\begin{snippetexample}{span-cifre-example}{Span di cifre}
    Per misurare lo span di cifre, si legge una lista di numeri casuali una sola
    volta e poi si chiede al partecipante di riprodurli nello stesso ordine. Per
    esempio: \texttt{8 1 7 3 4 9 4 2 8 5}. La maggior parte delle persone
    riesce a ricordare tra \(5\) e \(9\) elementi.
\end{snippetexample}

\begin{snippet}{numero-magico-miller}
    George Miller propose che \(7\), più o meno \(2\), fosse il ``numero
    magico'' che caratterizza molte prestazioni di memoria quando le persone
    devono ricordare liste di lettere, parole, numeri o altri elementi familiari
    e dotati di significato.
\end{snippet}

\begin{snippetdefinition}{ripetizione-definition}{Ripetizione}
    La \emph{ripetizione}, o \emph{rehearsal}, è un processo che mantiene attiva
    l'informazione nella memoria a breve termine attraverso la reiterazione.
\end{snippetdefinition}

\begin{snippetdefinition}{chunk-definition}{Chunk}
    Un \emph{chunk}, o \emph{unità di informazione}, è un'organizzazione di più
    bit di informazione in una singola unità dotata di significato.
\end{snippetdefinition}

\begin{snippetdefinition}{chunking-definition}{Chunking}
    Il \emph{chunking} è il processo di riorganizzazione degli elementi in unità
    significative, basato su somiglianza, significato o altre strutture già
    presenti nella memoria a lungo termine.
\end{snippetdefinition}

\begin{snippetexample}{chunking-example}{Chunking}
    La sequenza \texttt{1-9-8-4} può essere trattata come quattro cifre
    separate, ma se viene interpretata come l'anno \texttt{1984} forma una sola
    unità di informazione, lasciando più spazio disponibile nella memoria a
    breve termine.
\end{snippetexample}

\begin{snippetexample}{date-chunk-example}{Date come chunk}
    La sequenza \texttt{2001200620142021} può essere interpretata come \(16\)
    cifre indipendenti oppure come \(4\) anni: \texttt{2001}, \texttt{2006},
    \texttt{2014}, \texttt{2021}. Nel secondo caso è più facile memorizzare la
    sequenza corretta.
\end{snippetexample}

\begin{snippetdefinition}{memoria-lavoro-definition}{Memoria di lavoro}
    La \emph{memoria di lavoro} è una risorsa mnestica che permette di mantenere
    e manipolare temporaneamente le informazioni necessarie per ragionare,
    comprendere il linguaggio e risolvere problemi.
\end{snippetdefinition}

\begin{snippet}{componenti-memoria-lavoro}
    Secondo Alan Baddeley, la memoria di lavoro comprende:
    \begin{itemize}
        \item il \emph{circuito fonologico}, che mantiene e manipola
        informazioni verbali e acustiche;
        \item il \emph{taccuino visuospaziale}, che mantiene e manipola
        informazioni visive e spaziali;
        \item l'\emph{esecutivo centrale}, che controlla l'attenzione e
        coordina le altre componenti;
        \item il \emph{buffer episodico}, che combina informazioni della
        situazione attuale con informazioni recuperate dalla memoria a lungo
        termine.
    \end{itemize}
\end{snippet}

\end{document}
